\scanpage{49}
\mspage{47}
% archive sheet number: 48
On peut supposer si on veut $\unsure{U|X}$ affine, donc quasi-projectif.

\deleted{\unsure{En}} On admet \deleted{que} la résolution des singularités sur la clôture parfaite \deleted{de $K$} de $K$ en dim $\leq n$. Alors, quitte à passer à une extension radicale finie de $K$, on pourra trouver \inserted{un} morphisme surjectif \deleted{et} propre \inserted{birationnel}
\[
X'\longrightarrow X,
\]
avec $X'$ lisse. \deleted{\unsure{et}} La suite spectrale de descente de Deligne nous donne
\[
H^{\unsure{n|i}}(X_{\overline K},\unsure{\mathbb Q_\ell|Q_\ell})\Longleftarrow E^{pq}_2=H^p(\alpha\longmapsto H^q((X'/X)^{\alpha+1}_{\overline K},\mathbb Q_\ell).
\]
\deleted{Dans} \deleted{Cette suite spectrale}
\deleted{Nous supposons que l'on connaît le th. de monodromie pour $H^j$ pour $j<i$, et sans restriction sur la dimension de $X$.} L'hyp. de récurrence nous assure que \deleted{$E^{pq}_2$ \unsure{satisf}} l'action de $I$ sur $E^{pq}_2$ est quasi-unipotente quand $q<i$. \deleted{Pour $q=i$, on doit prendre} Parmi les termes $E^{pq}_2$ avec $p+q=i$, le seul dont on ne sache encore \deleted{\unsure{qu'il est}} si l'action de $I$ dessus est quasi unipotente est donc $E^{0,i}_2\subset H^i(X'_{\overline K},\mathbb Q_\ell)$.

\deleted{Mais} On trouve donc

\uline{Prop.} 9.6. \deleted{\inserted{\unsure{Supposons rés. des sing. sur} $K^{\unsure{p^{-\infty}|p^\infty}}$ \unsure{des schémas de tte dimension|des schémas de t.f. dim $\leq n$}.} \deleted{Pour} Supposons le th. de monodromie pour les $H^i(X_{\overline K},\mathbb Q_\ell)$ soit vrai pour $i<n$, \inserted{pour $X_K$ de t.f.,} \deleted{il faut et suffit qu'il soit vrai} \deleted{Alors il est vrai pour $\forall$} \deleted{\unsure{pour $\forall$ $X_K$}} et pour $i=n$, quand $X_K$ est lisse \deleted{\unsure{séparé de t.f.}} \inserted{quasi-projectif}, et \unsure{ceci étant|reste-t-il} après extension \unsure{pure|purement} inséparable $K'$ de $K$. \deleted{\unsure{et pour lisse}} \deleted{Alors} Alors il est vrai pour $H^n(X_{\overline K},\mathbb Q_\ell)$ pour les $X$ propres qui, sur la clôture \deleted{\unsure{inséparable}, \unsure{de}} $K^{p^{-\infty}}$, admet une résolution des singularités.}
